\documentclass[11pt]{article}

\usepackage[margin=0.95in]{geometry}
\usepackage{amsmath,booktabs,array,microtype,xcolor}
\usepackage[colorlinks=true,linkcolor=blue!55!black,
            citecolor=blue!55!black,urlcolor=blue!55!black]{hyperref}
\hypersetup{
  pdftitle={Ghosts, Geometry, and Reality in Fourth-Order Gravity},
  pdfauthor={GPT-5.6.sol; Claude Fable 5; Asger Alstrup Palm},
  pdfsubject={A thematic guide to the Weyl-gravity and reverse-foundations manuscript programme}
}

\newcommand{\keybox}[1]{%
  \begin{center}
  \fbox{\parbox{0.88\linewidth}{\centering #1}}
  \end{center}}
\newcommand{\paperlink}[2]{\href{#1}{\textbf{Paper #2}}}
\emergencystretch=1em

\title{Ghosts, Geometry, and Reality in Fourth-Order Gravity\\[4pt]
\large A thematic guide to the manuscript programme}
\author{GPT-5.6.sol\thanks{OpenAI model and principal programme author.}
\and Claude Fable 5\thanks{Anthropic model and computational coauthor on Paper 00.}
\and Asger Alstrup Palm\thanks{Programme orchestrator and corresponding human
contact; Honorary Professor, DTU Compute, Technical University of Denmark.
\texttt{asger@area9.dk}.}}
\date{17 August 2026}

\begin{document}
\maketitle

\begin{abstract}
This document is a guide, not another synthesis.  It explains how the papers
in the programme fit together, which question each paper is designed to
answer, which conclusions may be carried forward, and which stronger
interpretations are explicitly excluded.

The manuscript series follows four connected threads.  Papers 01--06 study
positive and indefinite completions of fourth-order systems and the
interaction problem.  Papers 07--13 move to gauge reduction, causal BV--BFV
transport, compact laboratories, observables, anomaly, and nonlinear
continuation.  Papers 14--18 study pure-Weyl black holes and the four-level
ghost classification.  Papers 19--22 separate physical assumptions from
mathematical carriers through Mannheim gravity, gauge algebra, reverse
foundations, and a Bateman--Turok counterfamily.  Papers 98 and 99 provide
physicist and general-audience entrances; Papers 90--92 and the computational
supplements expose selected interfaces and verification details.

The guide is organized by reading purpose rather than release order.  A
reader can enter through a physical question---ghosts, causal gauge theory,
black holes, galactic data, or mathematical foundations---and follow only the
papers and evidence rails needed for that question.
\end{abstract}

\section*{Authorship and use}

Asger Alstrup Palm, a computer scientist and Honorary Professor at DTU
Compute, orchestrates the programme and is the accountable human contact.
AI systems perform substantial derivation, programming, literature work,
adversarial review, and drafting.  Claude Fable~5 contributed to this guide.
Palm's affiliation identifies his academic role and does not imply
institutional endorsement.

The papers are research manuscripts, not substitutes for independent expert
review.  This guide describes their intended logical relationships; it does
not strengthen any theorem.  Follow each technical claim to its assumptions,
certificate, verifier, and explicit non-claims in the
\href{https://github.com/area9innovation/weyl-gravity}{public repository}.

\tableofcontents

\section{Choose an entrance}

The three public-facing documents have deliberately different jobs.

\begin{center}
\small
\begin{tabular}{>{\raggedright\arraybackslash}p{0.18\linewidth}
                >{\raggedright\arraybackslash}p{0.29\linewidth}
                >{\raggedright\arraybackslash}p{0.43\linewidth}}
\toprule
Reader & Begin with & Purpose\\
\midrule
General audience
& \href{99-how-to-build-a-universe.pdf}{\textbf{Paper 99}}
& Why a theory depends on physical, mathematical, logical, operational, and
empirical assumptions.\\
Physicist or referee
& \href{98-physicist-executive-summary.pdf}{\textbf{Paper 98}}
& Compact technical synthesis, exact claim boundaries, and routes into the
certificate atlas.\\
Reader choosing a specialist route
& \textbf{Paper 00} (this guide)
& Which manuscript answers which question, how papers depend on one another,
and where an inference must stop.\\
\bottomrule
\end{tabular}
\end{center}

For a single theme, the shortest routes are:
\begin{itemize}
\item \textbf{What is a ghost?}  Read Papers 01, 03, 04, 05, and 15.
\item \textbf{What survives gauge reduction?}  Read Papers 07, 08, 10, and 15.
\item \textbf{What survives interaction?}  Read Papers 05, 06, 11, and 13.
\item \textbf{What is the quantum obstruction?}  Read Paper 12, then Paper 98
for the Lorentzian boundary.
\item \textbf{What happens around Schwarzschild?}  Read Papers 16 and 17;
use Papers 14 and 18 for architecture and the separate static problem.
\item \textbf{What must Mannheim gravity assume?}  Read Papers 19 and 18.
\item \textbf{What changes when the mathematics changes?}  Read Papers 20,
21, and 22.
\end{itemize}

\section{How claims move through the series}

The series began from a four-level distinction that remains the safest way to
read it.  The word ``ghost'' can denote:
\begin{enumerate}
\item a local solution of a differential equation;
\item a classical direction surviving gauge symmetry, constraints, charges,
boundaries, and the symplectic radical;
\item a direction admitting consistent nonlinear continuation;
\item a state or particle in a positive quantum theory.
\end{enumerate}

A paper can establish one level while leaving the next untouched.  This is
why the guide gives every manuscript both a \emph{use} and a \emph{boundary}.
The boundary is part of the result, not a disclaimer added afterward.

Quantum claims also retain one or more exact dependency tags:
\begin{center}
\texttt{LOCAL-ALGEBRAIC}\quad
\texttt{EUCLIDEAN-SPECTRAL}\quad
\texttt{REDUCED-MODE}\quad
\texttt{LORENTZIAN-CAUSAL}.
\end{center}
A Euclidean determinant does not become a Lorentzian state.  A reduced mode
does not become a full metric/BV propagator.  A residual cohomology class does
not become a one-particle graviton without a separate transfer theorem.

\keybox{\textbf{How to use the guide.}
Choose a question, read the paper listed for that question, and stop at the
boundary stated in the same row.  Cross the boundary only through an explicit
paper, bridge note, or certificate.}

\section{Thread I: completion and interaction---Papers 01--06}

This thread starts with the Pais--Uhlenbeck oscillator, separates positive
metrics from vacuum and representation data, carries the distinction into
gravity, and then tests whether free completions survive interactions.

\begin{description}
\item[\paperlink{01-symplectic-diagonalization.pdf}{01}: canonical positive
symplectic diagonalization.]
Use this paper for its exact unequal-frequency normal form and positive
finite-dimensional metric.  Do not infer a field representation,
equal-frequency completion, interacting metric, or gravitational state.

\item[\paperlink{02-variational-fock.pdf}{02}: minimum distortion and field
representation.]
Use this paper to distinguish the variationally selected oscillator metric
from the ultraviolet requirements of a fourth-order field representation.
Do not treat a finite similarity transform as an automatically implementable
Fock-space transformation.

\item[\paperlink{03-fourth-order-vacuum.pdf}{03}: vacuum, Fock sectors, and
the Krein boundary.]
Use this paper for the classification of vacuum covariance, metric orbits,
Fock sectors, and the equal-frequency Jordan limit.  Do not infer that a
Krein carrier selects a positive physical state or survives interaction.

\item[\paperlink{04-fourth-order-gravity.pdf}{04}: gauge reduction in
fourth-order gravity.]
Use this paper to see how Lorentz covariance, helicity, constraints, and the
conformal Jordan boundary alter the scalar completion problem.  Do not model
the gravity system as independent copies of the oscillator or infer a
nonlinear quantum theory.

\item[\paperlink{05-interaction-obstructions.pdf}{05}: resonant interaction
obstruction.]
Use this paper for the scoped failure of the canonical analytic positive
metric under branch-changing interaction in fourth-order scalar systems.  It
does not exclude indefinite, nonlocal, singular, or nonperturbative
prescriptions.

\item[\paperlink{06-einstein-weyl-interaction-obstructions.pdf}{06}:
Einstein--Weyl interaction obstruction.]
Use this paper for cubic protection followed by a second-order metric failure
in a declared gravitational channel.  Do not promote the obstruction to all
interactions, all backgrounds, or strict pure-Weyl gravity.
\end{description}

\textbf{Thread logic:}
Paper 01 constructs a positive finite model; Papers 02--04 expose the extra
representation and gauge obligations; Papers 05--06 test the first nonlinear
promotion.  The correct conclusion is not ``positive'' or ``ghostly'' in the
abstract, but a typed map of which completion survives which gate.

\section{Thread II: causal and compact pure-Weyl theory---Papers 07--13}

This thread replaces local mode counting with full complexes, support
conditions, residual reduction, action-derived pairings, relational
observables, anomaly classes, and second-order continuation.

\begin{description}
\item[\paperlink{07-conformal-residual-cohomology-krein.pdf}{07}: residual
cohomology on the conformal cylinder.]
Use this paper for the declared zero-charge residual reduction and the
centered Weyl-square classes.  The classes are deformation or vertex classes,
not one-particle gravitons, and their positive residual pairing is not a
positive interacting Hilbert space.

\item[\paperlink{08-conformal-covariant-causal-transport.pdf}{08}: covariant
causal transport.]
Use this paper for the complete free BV--BFV transport between support
conditions and for the transported pairing.  Do not infer interacting causal
products, a positive Hadamard state, or a Lorentzian QME.

\item[\paperlink{09-relational-clocks-berger-d-cartan.pdf}{09}: a
backreacting phase clock.]
Use this paper for a controlled compact relational-clock construction and its
Cartan analysis.  The fixture is a model laboratory, not a realistic matter
clock, cosmology, or universal resolution of the problem of time.

\item[\paperlink{10-compact-einstein-maxwell-weyl-phase-space.pdf}{10}:
Einstein--Maxwell inside Weyl--Maxwell.]
Use this paper for the exact inclusion of the Einstein--Maxwell tangent and
the identification of additional nonradical branches on a compact product.
Classical quotient directions and Lee--Wald signs are not quantum particles
or probabilities.

\item[\paperlink{11-gravity-light-cyclic-causal-ell3.pdf}{11}: retained
gravity--light bracket.]
Use this paper for an exact mixed bracket surviving cyclic homological
reduction through the declared third operation.  Do not infer an all-order
interacting observable algebra or scattering construction.

\item[\paperlink{12-pure-weyl-one-loop-bv-anomaly.pdf}{12}: local one-loop BV
obstruction.]
Use this paper for the strict local Euclidean anomaly classification and the
formal compensator restoration in an enlarged theory.  The compensator
changes the field content; neither branch supplies a Lorentzian QME, positive
state, particles, or unitarity.

\item[\paperlink{13-compact-weyl-maxwell-second-order-tangent-cone.pdf}{13}:
formal second-order tangent cone.]
Use this paper for stabilizer moment maps, finite-harmonic continuation, and
declared bounded-resonance tests.  A formal second-order jet is not an exact
family, an all-orders solution, or a complete bounded nonlinear cone.
\end{description}

\textbf{Thread logic:}
Papers 07--08 construct the free residual and causal base.  Papers 09--11 ask
how clocks, phase space, and observables sit on that kind of structure.
Papers 12--13 test two different nonlinear promotions: quantum consistency
and classical second-order continuation.  Their outcomes do not substitute
for one another.

\section{Thread III: black holes and the four-level classification---Papers 14--18}

The black-hole thread separates operator factorization, endpoint selection,
resonance, causal response, and static charges.  It should not be read as one
undifferentiated ringdown claim.

\begin{description}
\item[\paperlink{14-pure-weyl-black-hole-radiation.pdf}{14}: black-hole
radiation architecture.]
Use this paper for the Ricci--Bach composition, horizon pairing, and the
operator architecture that motivates the later endpoint and resonance
theorems.  For theorem-strength endpoint and double-pole claims, continue to
Papers 16 and 17 rather than treating the architecture as their substitute.

\item[\paperlink{15-four-level-ghost-classification-phase1-synthesis.pdf}{15}:
four-level ghost classification.]
Use this paper to compare local solutions, classical reduced directions,
nonlinear continuation, and positive quantum states across the controlled
laboratories.  It is a synthesis of scoped results, not a universal rescue or
universal no-go theorem.

\item[\paperlink{16-lorentzian-endpoint-nonselection-pure-weyl.pdf}{16}:
Lorentzian endpoint non-selection.]
Use this paper for the theorem that future-horizon regularity and ordinary
one-sided radiative traces do not select the Einstein subsector of the
complete axial system in the stated scope.  Do not infer all-multipole,
polar, nonlinear, or quantum non-selection.

\item[\paperlink{17-pure-weyl-schwarzschild-extension-structure.pdf}{17}:
non-split extension and defective resonance.]
Use this paper for the repeated Regge--Wheeler self-extension, one certified
Smith type \((0,0,2)\), the nonzero double radial Green pole, and its retarded
mode-reduced parent.  It does not establish a full metric/BV causal resolvent,
global contour theorem, complete late-time expansion, real astrophysical
excitation, or detector waveform.

\item[\paperlink{18-static-bach-flat-black-hole-thermodynamics.pdf}{18}:
static Bach-flat charges and first laws.]
Use this paper for residual-basic normalization of Hamiltonian, entropy, and
signed temperature across the certified static Mannheim--Kazanas family.  It
is a static charge theorem, not a physical-process or radiative
thermodynamics theorem.
\end{description}

\textbf{Thread logic:}
Paper 14 supplies architecture; Paper 15 supplies the cross-laboratory claim
grammar; Paper 16 answers the endpoint-selection question; Paper 17 answers a
spectral and mode-reduced causal-response question; Paper 18 deliberately
returns to the separate static problem.

\section{Thread IV: assumptions and reverse foundations---Papers 19--22}

This thread asks which parts of a physical conclusion come from the equation,
which come from matter and boundary premises, which come from the gauge or
mathematical carrier, and which survive explicit counterexamples.

\begin{description}
\item[\paperlink{19-what-conformal-gravity-must-assume.pdf}{19}: what
conformal gravity must assume.]
Use this paper to separate the Mannheim--Kazanas radial solution from its
matter, conformal-frame, source-moment, boundary, and galactic-data premises.
Its NGC~3198 comparison is one frozen 39-point protocol, not population-level
model selection or a universal verdict on conformal gravity.

\item[\paperlink{20-is-the-gauge-algebra-an-assumption.pdf}{20}: direct check
of the Bach gauge algebra.]
Use this paper for a polynomial-jet search for Noether identities.  Within
the tested derivative orders and curvature degrees, it finds no additional
local generator after quotienting differential reparametrizations and
trivial symmetries.  It is a lower bound on completeness, not an all-orders
closure theorem.

\item[\paperlink{21-reverse-foundations-of-physics.pdf}{21}: reverse
foundations of physics.]
Use this paper for the literature synthesis and typed separation of logic,
existence theory, mathematical carrier, physical postulates, and one physical
obligation.  Its \(6\times6\times16=576\) atlas coordinates are questions,
not complete theories; populated cells and programme coverage do not compose
without interface theorems.

\item[\paperlink{22-bateman-turok-euclidean-torus-collapse.pdf}{22}: a
Bateman--Turok Green-tail counterfamily.]
Use this paper for the exact positive, nonseparable torus family that refutes
one deterministic all-field scaled-PL route.  It does not decide Gibbs
typicality, the interacting Witten or \(H^{-1}\) estimate, continuum
reconstruction, positivity, or any Lorentzian claim.
\end{description}

\textbf{Thread logic:}
Paper 19 separates a solution from its physical interpretation.  Paper 20
tests whether a familiar algebra is premise or consequence.  Paper 21 turns
that separation into a general research atlas.  Paper 22 demonstrates how the
atlas should behave when an attractive bridge fails: refute the bridge
precisely without inflating the result into a verdict on the entire programme.

\section{Public syntheses, supplements, and bridge notes}

\subsection{Public syntheses}

\begin{description}
\item[\href{98-physicist-executive-summary.pdf}{\textbf{Paper 98: Reverse
physics and pure-Weyl gravity.}}]
The physicist-facing synthesis.  Use it for the compact programme argument,
exact dependency boundaries, principal numerical comparisons, and reviewer
route into the interactive evidence atlas.

\item[\href{99-how-to-build-a-universe.pdf}{\textbf{Paper 99: How to Build a
Universe.}}]
The general-audience introduction.  Use it for the theory-as-a-stack idea,
the purpose of the 576-coordinate atlas, and the role of exposition and
verification in AI-assisted research.
\end{description}

\subsection{Computational supplements and archival manuscript}

The supplements document derivations, conventions, tables, and reproduction
commands without changing the claim boundary of their parent paper:
\begin{itemize}
\item \href{07-08-conformal-residual-cohomology-computational-supplement.pdf}{Papers
07--08 computational supplement};
\item \href{09-relational-clocks-berger-d-cartan-computational-supplement.pdf}{Paper
09 computational supplement};
\item \href{12-pure-weyl-one-loop-bv-anomaly-computational-supplement.pdf}{Paper
12 computational supplement};
\item \href{16-endpoint-nonselection-computational-supplement.pdf}{Paper 16
computational supplement}.
\end{itemize}
The \href{07-08-conformal-residual-cohomology-archive.pdf}{combined Papers
07--08 archival manuscript} preserves the earlier paired presentation; use
the separate Papers 07 and 08 for the cleaner residual/causal division.

\subsection{Bridge notes 90--92}

Bridge notes translate between communities and expose interfaces; they are
not additional headline-paper conclusions.
\begin{description}
\item[\href{90-cyclic-green-transfer-bridge.pdf}{\textbf{Paper 90: cyclic
Green transfer.}}]
Connects Green-hyperbolic complexes, support-local cyclic deformation
retracts, and conformal detour constructions.  Use it with Papers 08 and 11.

\item[\href{91-charge-fibre-taub-bridge.pdf}{\textbf{Paper 91: charge-fibre
Taub bridge.}}]
Connects compact pure-extra obstruction and a balanced Einstein--extra
second-order extension.  Use it with Papers 10 and 13; its formal second-order
jet is not an all-orders family.

\item[\href{92-extra-axial-lee-wald-bridge.pdf}{\textbf{Paper 92: compact
currents and black-hole endpoints.}}]
Compares extra classical Lee--Wald directions with adjacent positive-metric,
BRST/Fock, and boundary-selection questions.  Use it between Papers 10,
16, and 17; classical signs remain distinct from quantum norms.
\end{description}

\section{Reading routes by research question}

\begin{center}
\small
\begin{tabular}{>{\raggedright\arraybackslash}p{0.31\linewidth}
                >{\raggedright\arraybackslash}p{0.59\linewidth}}
\toprule
Question & Route\\
\midrule
Can a fourth-order free theory have a positive completion?
& 01 \(\rightarrow\) 02 \(\rightarrow\) 03 \(\rightarrow\) 04.\\
Does that completion survive interaction?
& 05 \(\rightarrow\) 06; compare 13 for a different nonlinear gate.\\
What survives the classical gauge quotient?
& 07 \(\rightarrow\) 08 \(\rightarrow\) 10; synthesize with 15.\\
Can an observable or clock be transported?
& 09 and 11, with Bridge 90 for the causal-transfer interface.\\
Where does quantization first fail or remain incomplete?
& 12 \(\rightarrow\) 98; keep Euclidean anomaly and Lorentzian pseudo-state
as separate questions.\\
Do black-hole boundary conditions remove the extra sector?
& 14 \(\rightarrow\) 16 \(\rightarrow\) 17; use 18 only for the separate
static charge problem.\\
Does conformal gravity explain rotation curves?
& 19, then 18 for static geometry and 21 for the assumption map.\\
Is the gauge algebra assumed or derived?
& 20 \(\rightarrow\) 21.\\
What if Choice, Hilbert space, or the continuum is changed?
& 21, using 03 for Krein structure and 22 for a concrete Euclidean
counterexample.\\
How should an external reviewer begin?
& 98 \(\rightarrow\) relevant technical paper \(\rightarrow\) certificate
\(\rightarrow\) independent verifier.\\
\bottomrule
\end{tabular}
\end{center}

\section{How to inspect a claim}

The \href{../foundations/site/index.html}{Reverse Physics Atlas} is the
interactive companion to this guide.  It exposes the dimensions, matrix
cells, theory profiles, programme assemblies, end-to-end journeys, Weyl BV
routes, literature records, certificates, and formal-proof passports.

For a computer-assisted result, follow
\[
\begin{aligned}
 \text{paper claim}
 &\longrightarrow \text{producer}
 \longrightarrow \text{machine-readable certificate}\\
 &\longrightarrow \text{independent verifier}
 \longrightarrow \text{adversarial or mutation test}.
\end{aligned}
\]
Re-running the producer is reproduction, not independent verification.
Exact algebra and numerical evidence are different types.  A skip, missing
input, or timeout is not a pass.

The manuscript sequence alone is not a dependency proof.  Two papers may use
different actions, backgrounds, carriers, support conditions, or empirical
protocols.  A later paper can sharpen an earlier question without promoting
all of its claims.  The certificate and non-claim boundary determine what may
be transferred.

\section{What the whole series does and does not claim}

Across the four threads, the programme establishes examples of positive and
Krein free completions, interaction obstructions, residual and causal gauge
constructions, additional nonradical classical directions, nonlinear balance
and obstruction, a strict local Euclidean anomaly, a full-complex indefinite
BRST--Hadamard pseudo-state pair, black-hole endpoint non-selection and a
defective mode-reduced response, bounded empirical controls, and a
deterministic Euclidean counterfamily.

Those ingredients do not compose automatically into a complete theory.  The
series does not establish a positive full-BV particle space, Lorentzian
interacting quantum gravity, scattering or unitarity, a complete causal
black-hole waveform, population-level galactic model selection, a continuum
Bateman--Turok reconstruction, or a universal weakest-foundation theorem.

\keybox{\textbf{Role of Paper 00.}
Paper 99 explains why the programme matters.  Paper 98 states the compact
physicist-level synthesis.  Paper 00 tells the reader where to go next and
where each route must stop.}

\end{document}
